
\begin{figure*}[t!]
\centering
\resizebox{\textwidth}{!}{%
\begin{tikzpicture}[x=0.80cm,y=0.80cm,
    >=Latex,
    line cap=round,line join=round,
    every node/.style={inner sep=0pt},
    stage/.style={font=\sffamily\fontsize{8.5}{10}\selectfont\bfseries},
    textlabel/.style={font=\sffamily\small},
    note/.style={font=\sffamily\fontsize{7.5}{9}\selectfont,text=ddgray}]

% Colors follow the paper-style blue / red / purple / green palette.
\definecolor{ddblue}{RGB}{91,132,180}
\definecolor{ddbluefill}{RGB}{217,230,244}
\definecolor{ddred}{RGB}{199,106,100}
\definecolor{ddredfill}{RGB}{248,222,218}
\definecolor{ddpurple}{RGB}{169,154,187}
\definecolor{ddpurplefill}{RGB}{250,247,252}
\definecolor{ddgreen}{RGB}{102,142,120}
\definecolor{ddgreenfill}{RGB}{247,250,243}
\definecolor{ddgray}{RGB}{91,101,111}

% Column headings.
\node[stage] at (5.47,7.02) {Shared augmentation};
\node[note] at (5.47,6.62) {within each view};
\node[stage] at (9.52,7.02) {Frozen encoder};
\node[font=\footnotesize,text=ddgray] at (9.52,6.62) {$f_{\theta_{\mathcal{T}}}$};

% Real and synthetic image stacks.
\node[textlabel,text=ddblue] at (1.77,6.62) {Real batch};
\foreach \xx/\yy in {1.43/5.19,1.32/5.08,1.21/4.97,1.10/4.86}{%
  \draw[fill=ddbluefill,draw=ddblue,line width=.70pt]
    (\xx,\yy) rectangle ++(1.02,1.01);
}
\node[font=\small,text=ddblue] at (1.61,5.365) {$\mathcal{B}$};
\node[textlabel,text=ddred] at (1.77,3.91) {Synthetic set};

\foreach \xx/\yy in {1.43/2.48,1.32/2.37,1.21/2.26,1.10/2.15}{%
  \draw[fill=ddredfill,draw=ddred,line width=.70pt]
    (\xx,\yy) rectangle ++(1.02,1.01);
}
\node[font=\small,text=ddred] at (1.61,2.655) {$\mathcal{S}$};

% Branch arrows: the same real and synthetic images feed both views.
\draw[->,draw=ddblue,line width=.95pt] (2.48,5.26)--(4.23,5.26);
\draw[draw=ddblue,line width=.95pt] (3.31,5.26)--(3.31,2.48);
\draw[->,draw=ddblue,line width=.95pt] (3.31,2.48)--(4.23,2.48);
\fill[ddblue] (3.31,5.26) circle (.050);

\draw[draw=ddred,line width=.95pt] (2.51,2.66)--(3.68,2.66);
\draw[draw=ddred,line width=.95pt] (3.68,1.87)--(3.68,4.62);
\draw[->,draw=ddred,line width=.95pt] (3.68,4.62)--(4.23,4.62);
\draw[->,draw=ddred,line width=.95pt] (3.68,1.87)--(4.23,1.87);
\fill[ddred] (3.68,2.66) circle (.050);

% View 1 and view 2 augmentation boxes.
\foreach \by/\v in {4.19/1,1.40/2}{%
  \draw[rounded corners=.17cm,draw=ddpurple,line width=.8pt,fill=ddpurplefill]
    (4.27,\by) rectangle (6.68,\by+2.07);
  \node[textlabel,text=ddgray] at (5.47,\by+1.73)
    {View \v: $a^{(\v)}$};
  \draw[rounded corners=.055cm,draw=ddblue,line width=.55pt,fill=ddbluefill]
    (4.53,\by+.82) rectangle (6.40,\by+1.36);
  \node[font=\footnotesize,text=ddblue] at (5.47,\by+1.09)
    {$a^{(\v)}(\mathcal{B})$};
  \draw[rounded corners=.055cm,draw=ddred,line width=.55pt,fill=ddredfill]
    (4.53,\by+.17) rectangle (6.40,\by+.71);
  \node[font=\footnotesize,text=ddred] at (5.47,\by+.44)
    {$a^{(\v)}(\mathcal{S})$};
}
\node[note] at (5.47,1.00) {Independent views};

% Arrows from the augmented inputs to the frozen feature network.
\foreach \yy in {5.26,2.48}{\draw[->,ddblue,line width=.95pt] (6.68,\yy)--(7.79,\yy);}
\foreach \yy in {4.62,1.87}{\draw[->,ddred,line width=.95pt] (6.68,\yy)--(7.79,\yy);}

% Shared encoder enclosure.
\begin{scope}[xshift=1.2mm]
\draw[rounded corners=.16cm,draw=ddgray!70,fill=black!2,line width=.75pt]
  (7.64,1.38) rectangle (11.10,6.26);

% Two views of the same fixed network; all connections and nodes editable.
\foreach \cy in {5.06,2.26}{%
  \foreach \ya in {-.43,0,.43}{%
    \foreach \yb in {-.63,-.21,.21,.63}{%
      \draw[ddgray!85,line width=.27pt] (7.98,\cy+\ya)--(8.49,\cy+\yb);
    }
  }
  \foreach \yb in {-.63,-.21,.21,.63}{%
    \foreach \yc in {-.29,.29}{%
      \draw[ddgray!85,line width=.27pt] (8.49,\cy+\yb)--(9.05,\cy+\yc);
    }
  }
  \foreach \ya in {-.43,0,.43}{\draw[fill=black!3,draw=ddgray,line width=.55pt] (7.98,\cy+\ya) circle (.11);}
  \foreach \yb in {-.63,-.21,.21,.63}{\draw[fill=black!3,draw=ddgray,line width=.55pt] (8.49,\cy+\yb) circle (.11);}
  \foreach \yc in {-.29,.29}{\draw[fill=black!3,draw=ddgray,line width=.55pt] (9.05,\cy+\yc) circle (.11);}
  \draw[->,draw=ddgray,line width=.7pt] (9.24,\cy)--(9.62,\cy);
  % Upper (real) feature bar.
  \foreach \k/\shade in {0/20,1/55,2/16,3/44,4/68}{%
    \fill[ddblue!\shade!white] (9.74+\k*.24,\cy+.07) rectangle ++(.235,.30);
  }
  \draw[draw=ddblue,line width=.55pt] (9.74,\cy+.07) rectangle (10.94,\cy+.37);
  % Lower (synthetic) feature bar.
  \foreach \k/\shade in {0/55,1/13,2/29,3/10,4/41}{%
    \fill[ddred!\shade!white] (9.74+\k*.24,\cy-.57) rectangle ++(.235,.30);
  }
  \draw[draw=ddred,line width=.55pt] (9.74,\cy-.57) rectangle (10.94,\cy-.27);
}
\node[note] at (9.37,3.79) {shared weights};
\node[note,font=\sffamily\footnotesize] at (9.37,1.00) {$\ell_2$-normalized features};
\end{scope}

% Stage headings share a baseline and align with their visual blocks.
\node[stage] at (13.92,7.02) {Prototype assignments};
\node[note] at (13.92,6.62) {cosine similarity + softmax};
\node[stage] at (19.33,7.02) {Joint distribution};
\node[stage] at (22.95,7.02) {Information objective};

% Encoder outputs enter two sparse, bipartite assignment diagrams.
\draw[->,ddgray,line width=.9pt] (11.26,5.08)--(12.11,5.08);
\draw[->,ddgray,line width=.9pt] (11.26,2.29)--(12.11,2.29);

% Instead of Q heatmaps, the two views use bipartite graphs.
% Blue circles are real examples; red squares are synthetic prototypes.
% Edge darkness indicates an illustrative soft assignment weight.
\begin{scope}[xshift=2.4mm]
\foreach \v/\by in {1/4.13,2/1.35}{%
  \draw[rounded corners=.15cm,fill=ddpurplefill,draw=ddpurple!80,
       line width=.70pt] (11.94,\by+.15) rectangle (15.29,\by+1.98);
  \foreach \ra in {.48,1.05,1.61}{%
    \foreach \pr in {.46,.86,1.27,1.67}{%
      \draw[ddpurple!26,line width=.34pt]
        (12.35,\by+\ra)--(14.85,\by+\pr);
    }
  }
  % Illustrative soft assignments differ in both strength and preferred prototype.
  \ifnum\v=1
    \draw[ddpurple!95,line width=1.58pt] (12.35,\by+1.61)--(14.85,\by+1.67);
    \draw[ddpurple!82,line width=1.42pt] (12.35,\by+1.05)--(14.85,\by+.86);
    \draw[ddpurple!86,line width=1.48pt] (12.35,\by+.48)--(14.85,\by+.46);
    \draw[ddpurple!57,line width=.84pt] (12.35,\by+1.05)--(14.85,\by+1.27);
    \draw[ddpurple!52,line width=.76pt] (12.35,\by+.48)--(14.85,\by+.86);
  \else
    \draw[ddpurple!88,line width=1.46pt] (12.35,\by+1.61)--(14.85,\by+1.27);
    \draw[ddpurple!95,line width=1.58pt] (12.35,\by+1.05)--(14.85,\by+1.67);
    \draw[ddpurple!82,line width=1.40pt] (12.35,\by+.48)--(14.85,\by+.86);
    \draw[ddpurple!61,line width=.90pt] (12.35,\by+1.61)--(14.85,\by+1.67);
    \draw[ddpurple!48,line width=.70pt] (12.35,\by+.48)--(14.85,\by+.46);
  \fi
  % The nodes sit above the links for a clean outline.
  \foreach \ra in {.48,1.05,1.61}{%
     \draw[fill=ddbluefill,draw=ddblue,line width=.77pt]
          (12.35,\by+\ra) circle (.115);
  }
  \foreach \pr in {.46,.86,1.27,1.67}{%
     \draw[fill=ddredfill,draw=ddred,line width=.78pt,rounded corners=.018cm]
          (14.74,\by+\pr-.105) rectangle (14.95,\by+\pr+.105);
  }
  \node[font=\footnotesize] at (13.62,\by-.20)
       {$\mathbf{Q}_{\v}\in\mathbb{R}^{B\times m}$};
}
\end{scope}

% All four right-side arrows span 1.00 unit horizontally.
\draw[->,ddgray,line width=.90pt] (15.59,5.17) -- (16.59,4.1034);
\draw[->,ddgray,line width=.90pt] (15.59,2.39) -- (16.59,3.5366);
\fill[white] (16.90,3.82) circle (.42);
\draw[draw=ddpurple,line width=.94pt] (16.90,3.82) circle (.42);
\node[font=\large,text=ddpurple!90!black] at (16.90,3.82) {$\times$};
\draw[->,ddgray,line width=.90pt] (17.32,3.82) -- (18.32,3.82);

% Center the joint matrix on the connecting arrows.
% Symmetric colors make the structure of P visually explicit.
\foreach \rowidx in {0,1,2,3}{%
  \foreach \colidx in {0,1,2,3}{%
    \pgfmathtruncatemacro{\sep}{abs(\rowidx+\colidx-3)}
    \pgfmathtruncatemacro{\shade}{max(12,84-29*\sep-2*abs(\rowidx-\colidx))}
    \fill[ddgreen!\shade!white]
       (18.37+\colidx*.48,2.86+\rowidx*.48) rectangle ++(.48,.48);
    \draw[white,line width=.45pt]
       (18.37+\colidx*.48,2.86+\rowidx*.48) rectangle ++(.48,.48);
  }
}
\draw[draw=ddgreen!82,line width=.88pt] (18.37,2.86) rectangle (20.29,4.78);
\node[font=\small] at (19.33,2.29) {$J=\frac{1}{B}\mathbf{Q}_1^\top\mathbf{Q}_2$};
\node[font=\small] at (19.33,1.50) {$P=\frac{1}{2}(J+J^\top)$};
\draw[->,ddgray,line width=.90pt] (20.34,3.82) -- (21.34,3.82);
  
\node[stage] at (22.80,4.12) {Maximize};
\node[font=\small] at (22.80,3.51) {$I_P(C_1;C_2)$};

% Keep the existing feedback: only synthetic pixels are updated.
\draw[draw=ddred,dashed,line width=1.00pt,->]
  (22.80,2.86)--(22.80,.52)--(1.75,.52)--(1.75,2);
\node[note,text=ddred,fill=white,inner xsep=5pt] at (12.28,.52)
  {Update synthetic pixels};

\end{tikzpicture}%
}
\caption{Overview of the proposed self-supervised dataset distillation method.
\textbf{MIPA (Mutual Information Prototype Alignment)}}
\label{fig:ssl-dd-pipeline}
\end{figure*}
